\subsection[Identité autoduale de Planck sur la périodicité d'ordre 108]{Identité autoduale de
Planck dans la réalisation circulaire de l'endomorphisme temporel d'ordre 108}
\label{sec:identidad-autodual-planck-ct108-rev3}
\index{Planck!identité autoduale}
\index{calendrier d’ordre 108!réalisation circulaire}
\index{Compton!lectures réduite et complète}

La droite d’action construite dans la section précédente admet une
réalisation circulaire compatible avec la périodicité temporelle d’ordre 108.
Cette réalisation coordonne trois objets qui possèdent déjà leurs types propres : la
section d’action, la géométrie d’un tour et la transduction
gravitationnelle. La coordination qui en résulte fournit une identité de
Planck interne et prépare la droite positive sur laquelle agira ensuite
l’opérateur général de masse.

\subsubsection{Cartes orientées d’action et réalisation circulaire déclarée}

Soient
\[
 \mathcal L_S^+,\quad \mathcal L_T^+,\quad \mathcal L_C^+,
 \quad \mathcal L_L^+,\quad \mathcal L_M^+,\quad \mathcal L_G^+
\]
les demi-droites positives d’action, de temps, de vitesse, de longueur et de masse.
Les coordonnées antérieure et postérieure au retour s’écrivent
\begin{equation}
 \hbar_{\rm pre}=H_5\,10^{-34}\mathcal U_S,
 \qquad
 \hbar_{\rm ret}=\eta_{\rm ret}\,10^{-34}\mathcal U_S,
 \qquad
 R_{\rm act}=\frac{\hbar_{\rm pre}}{\hbar_{\rm ret}}
 =\frac{H_5}{\eta_{\rm ret}}.
 \label{eq:planck-ct108-cartas-pre-ret-rev3}
\end{equation}
L’indice \(a\in\{\mathrm{pre},\mathrm{ret}\}\) désignera l’une quelconque
des deux cartes. Dans chacune, la distinction exacte est conservée
\begin{equation}
 h_a=2\pi\hbar_a .
 \label{eq:planck-ct108-h-hbar-rev3}
\end{equation}

Fixons une durée élémentaire \(t_0\in\mathcal L_T^+\), une vitesse
interne \(c_{\rm int}\in\mathcal L_C^+\) et
\(\ell_0=c_{\rm int}t_0\). La réalisation circulaire déclarée de la
périodicité temporelle d’ordre 108 est déterminée par
\begin{equation}
 T_{108}=108t_0,
 \qquad
 \omega_{108}=\frac{2\pi}{T_{108}},
 \qquad
 R_{108}=\frac{c_{\rm int}}{\omega_{108}}
 =\frac{54}{\pi}\ell_0,
 \label{eq:planck-ct108-realizacion-circular-rev3}
\end{equation}
et par la longueur du tour complet
\begin{equation}
 C_{108}=2\pi R_{108}=c_{\rm int}T_{108}=108\ell_0.
 \label{eq:planck-ct108-perimetro-rev3}
\end{equation}
Ainsi, \(R_{108}\) est le rayon de la réalisation et \(C_{108}\) son
périmètre. La fréquence angulaire, la fréquence cyclique et les deux
constantes d’action sont reliées par
\(\omega_{108}=2\pi/T_{108}\) et \(h_a=2\pi\hbar_a\).

\subsubsection{Correspondance entre la périodicité d’ordre 108 et les lectures de Compton}

Dans la carte \(a\), le quantum circulaire et sa section de masse sont
\begin{equation}
 E_{108,a}=\hbar_a\omega_{108},
 \qquad
 m_{108,a}=\frac{E_{108,a}}{c_{\rm int}^{2}}
 =\frac{\hbar_a\omega_{108}}{c_{\rm int}^{2}}.
 \label{eq:planck-ct108-cuanto-circular-rev3}
\end{equation}

\begin{theorem}[Identité entre la périodicité circulaire et les lectures de Compton]
\label{thm:planck-ct108-compton-rev3}
Pour chaque carte \(a\in\{\mathrm{pre},\mathrm{ret}\}\), les lectures de
Compton réduite et complète satisfont
\begin{equation}
 \boxed{
 \bar\lambda_{\rm C}(m_{108,a})
 =\frac{\hbar_a}{m_{108,a}c_{\rm int}}=R_{108},
 \qquad
 \lambda_{\rm C}(m_{108,a})
 =\frac{h_a}{m_{108,a}c_{\rm int}}=C_{108}.}
 \label{eq:planck-ct108-identidad-compton-rev3}
\end{equation}
\end{theorem}

\begin{proof}
La substitution de
\eqref{eq:planck-ct108-cuanto-circular-rev3} produit
\[
 \frac{\hbar_a}{m_{108,a}c_{\rm int}}
 =\frac{c_{\rm int}}{\omega_{108}}=R_{108}.
\]
La deuxième égalité résulte de l’application simultanée de
\eqref{eq:planck-ct108-h-hbar-rev3} et de
\eqref{eq:planck-ct108-perimetro-rev3}.
\end{proof}

Le théorème est homogène en
\((\hbar_a,c_{\rm int},t_0)\) : la même périodicité détermine le rayon,
le périmètre et les deux cartes de Compton. La réalisation circulaire est la
primitive géométrique déclarée ; les égalités précédentes sont ses
conséquences exactes.

\subsubsection{Sélecteur circulaire et unité interne de Planck}

La droite gravitationnelle est coordonnée à la carte \(a\) par le transducteur
homogène
\begin{equation}
 G_a(\theta)=\theta\,
 \frac{c_{\rm int}^{5}t_0^2}{\hbar_a}
 \in\mathcal L_G^+,
 \qquad \theta\in\mathbb R_{>0}.
 \label{eq:planck-ct108-transductor-rev3}
\end{equation}
Son rayon gravitationnel réduit, évalué sur le quantum circulaire, est
\begin{equation}
 r_{{\rm g},a}(\theta)
 :=\frac{G_a(\theta)m_{108,a}}{c_{\rm int}^{2}}
 =\theta\frac{\pi}{54}\ell_0.
 \label{eq:planck-ct108-radio-gravitatorio-rev3}
\end{equation}

\begin{proposition}[Sélecteur circulaire de Planck]
\label{prop:selector-circular-planck-ct108-rev3}
Au sein de la réalisation circulaire
\eqref{eq:planck-ct108-realizacion-circular-rev3}, la condition
\(r_{{\rm g},a}(\theta)=R_{108}\) sélectionne une unique coordonnée
positive,
\begin{equation}
 \boxed{\theta_{108}^{\rm circ}=\left(\frac{54}{\pi}\right)^2.}
 \label{eq:theta-circular-planck-ct108-rev3}
\end{equation}
La coordonnée est indépendante de la carte orientée \(a\).
\end{proposition}

\begin{proof}
Les équations
\eqref{eq:planck-ct108-realizacion-circular-rev3} et
\eqref{eq:planck-ct108-radio-gravitatorio-rev3} réduisent la condition à
\[
 \theta\frac{\pi}{54}\ell_0=\frac{54}{\pi}\ell_0.
\]
La positivité de \(\ell_0\) et la stricte monotonie de l’application
\(\theta\mapsto r_{{\rm g},a}(\theta)\) fournissent la solution et son
unicité.
\end{proof}

Soit \(G_{108,a}^{\rm circ}:=G_a(\theta_{108}^{\rm circ})\). Les unités internes
associées à cette carte satisfont
\begin{align}
 \ell_{{\rm P},a}^{\rm circ}
 &:=\sqrt{\frac{\hbar_a G_{108,a}^{\circ}}{c_{\rm int}^{3}}}
 =R_{108},
 &
 t_{{\rm P},a}^{\circ}
 &:=\sqrt{\frac{\hbar_a G_{108,a}^{\circ}}{c_{\rm int}^{5}}}
 =\omega_{108}^{-1},
 \label{eq:planck-ct108-longitud-tiempo-rev3}\\
 m_{{\rm P},a}^{\circ}
 &:=\sqrt{\frac{\hbar_a c_{\rm int}}{G_{108,a}^{\circ}}}
 =m_{108,a},
 &
 E_{{\rm P},a}^{\circ}
 &:=m_{{\rm P},a}^{\circ}c_{\rm int}^{2}=E_{108,a}.
 \label{eq:planck-ct108-masa-energia-rev3}
\end{align}
L’égalité des rayons est donc publiée comme
\begin{equation}
 R_{108}=\bar\lambda_{\rm C}=r_{\rm g}=\ell_{\rm P}^{\rm circ},
 \qquad
 C_{108}=\lambda_{\rm C}=2\pi r_{\rm g}=2\pi\ell_{\rm P}^{\rm circ}.
 \label{eq:planck-ct108-tres-lectores-rev3}
\end{equation}

La convention de ce théorème est le rayon gravitationnel réduit
\(r_{\rm g}=Gm/c_{\rm int}^{2}\). Le rayon de Schwarzschild satisfait
\(r_{\rm s}=2r_{\rm g}\) ; si cette seconde convention est adoptée, son croisement avec
\(\bar\lambda_{\rm C}\) se produit en
\(m=m_{{\rm P},a}^{\circ}/\sqrt2\). La déclaration du rayon fait donc
partie de la carte et évite de transférer par inadvertance un facteur deux
entre les deux identités.

La coordonnée \(\theta_{108}^{\rm circ}\) appartient à la spécialisation
circulaire déclarée. Le sélecteur gravitationnel général, formulé sur
l’holonomie enrichie et sa réponse spectrale, conserve sa résidence dans
la \cref{sec:selector-gravitatorio-rev11}. Ainsi, l’identité
circulaire fixe une normalisation interne de Planck et le critère général
fixe les conditions que doit satisfaire la réalisation physique de la
section gravitationnelle.

\subsubsection{Covariance pré-retour/post-retour}

Le produit des sections réciproques d’action et de gravité conserve la
géométrie circulaire :
\begin{equation}
 \frac{m_{108,\rm pre}}{m_{108,\rm ret}}=R_{\rm act},
 \qquad
 \frac{G_{108,\rm pre}^{\circ}}{G_{108,\rm ret}^{\circ}}
 =R_{\rm act}^{-1},
 \qquad
 \hbar_a G_{108,a}^{\circ}
 =\theta_{108}^{\rm circ}c_{\rm int}^{5}t_0^2.
 \label{eq:planck-ct108-covariancia-pre-ret-rev3}
\end{equation}
En particulier, \(\ell_{{\rm P},\rm pre}^{\rm circ}\) et
\(\ell_{{\rm P},\rm ret}^{\rm circ}\) coïncident avec \(R_{108}\). La mémoire
bidirectionnelle réside dans le quotient des masses et dans son réciproque
gravitationnel ; l’échelle géométrique réside dans le produit invariant.

La réciprocité admet une paramétrisation symétrique qui sépare l’échelle
commune de l’orientation du retour. Définissons
\begin{equation}
 \begin{aligned}
 \hbar_{\rm geom}
 &:=\sqrt{\hbar_{\rm pre}\hbar_{\rm ret}},
 &
 G_{108,\rm geom}^{\circ}
 &:=\sqrt{G_{108,\rm pre}^{\circ}G_{108,\rm ret}^{\circ}},
 &
 \xi_{\rm act}&:=\frac12\log R_{\rm act}.
 \end{aligned}
 \label{eq:planck-ct108-parametrizacion-simetrica-rev3}
\end{equation}
Alors
\[
 \hbar_{\rm pre}=\hbar_{\rm geom}e^{\xi_{\rm act}},
 \qquad
 \hbar_{\rm ret}=\hbar_{\rm geom}e^{-\xi_{\rm act}},
\]
tandis que les sections gravitationnelles conjuguées satisfont
\[
 G_{108,\rm pre}^{\circ}
 =G_{108,\rm geom}^{\circ}e^{-\xi_{\rm act}},
 \qquad
 G_{108,\rm ret}^{\circ}
 =G_{108,\rm geom}^{\circ}e^{\xi_{\rm act}}.
\]
La coordonnée \(\xi_{\rm act}\) conserve l’orientation
pré-retour/post-retour ; le produit \(\hbar_aG_{108,a}^{\circ}\)
conserve l’échelle géométrique. Cette paramétrisation des cartes d’action
n’introduit pas à elle seule des vitesses de propagation distinctes.

\subsubsection{Droite autoduale déployée des masses}

Une carte \(a\) étant fixée, soit l’algèbre déployée
\[
 \mathbb D_{\rm sp}:=\mathbb R[j]/(j^2-1),
 \qquad
 P_\pm:=\frac{I\pm j}{2},
\]
et considérons sa décomposition spectrale
\[
 \mathbb D_{\rm sp}
 =\mathbb R P_+\oplus\mathbb R P_-,
 \qquad
 P_\pm^2=P_\pm,\quad
 P_+P_-=0,\quad
 P_++P_-=I.
\]
La conjugaison déployée échange les projecteurs,
\[
 \overline{uP_++vP_-}=vP_++uP_-,
\]
et sa norme est
\[
 N_{\rm sp}(uP_++vP_-):=uv.
\]

Pour \(m\in\mathcal L_M^+\), définissons la section
\begin{equation}
 Z_a(m)
 :=\frac{m_{{\rm P},a}^{\circ}}{m}P_+
   +\frac{m}{m_{{\rm P},a}^{\circ}}P_- .
 \label{eq:planck-ct108-seccion-partida-masa-rev3}
\end{equation}
L’unité \(m_{{\rm P},a}^{\circ}\) détermine également l’involution
\begin{equation}
 \iota_{{\rm P},a}(m)
 =\frac{(m_{{\rm P},a}^{\circ})^2}{m}.
 \label{eq:planck-ct108-involucion-masas-rev3}
\end{equation}

\begin{theorem}[Droite autoduale déployée des masses]
\label{thm:planck-ct108-recta-partida-masa-rev3}
Pour tout \(m\in\mathcal L_M^+\),
\begin{equation}
 \boxed{
 N_{\rm sp}(Z_a(m))=1,
 \qquad
 \overline{Z_a(m)}
 =Z_a\!\left(\iota_{{\rm P},a}(m)\right).}
 \label{eq:planck-ct108-norma-conjugacion-rev3}
\end{equation}
La conjugaison possède un unique point fixe positif,
\[
 m=m_{{\rm P},a}^{\circ},
 \qquad
 Z_a(m_{{\rm P},a}^{\circ})=I.
\]
\end{theorem}

\begin{proof}
La norme est le produit des deux coordonnées réciproques :
\[
 \frac{m_{{\rm P},a}^{\circ}}m
 \frac m{m_{{\rm P},a}^{\circ}}=1.
\]
La conjugaison échange les deux coordonnées, ce qui équivaut à remplacer
\(m\) par \((m_{{\rm P},a}^{\circ})^2/m\). Le point fixe positif satisfait
\(m^2=(m_{{\rm P},a}^{\circ})^2\).
\end{proof}

Avec la rapidité
\begin{equation}
 x_a(m)=\log\frac{m}{m_{{\rm P},a}^{\circ}},
 \label{eq:planck-ct108-rapidez-masa-rev3}
\end{equation}
la section et sa conjuguée s’écrivent
\[
 Z_a(m)=e^{-x_a}P_++e^{x_a}P_-,
 \qquad
 \overline{Z_a(m)}
 =e^{x_a}P_++e^{-x_a}P_-.
\]
Leurs lecteurs pair et impair sont
\begin{equation}
 Q_a(m):=\frac{e^{x_a}+e^{-x_a}}2=\cosh x_a,
 \qquad
 A_a(m):=\frac{e^{x_a}-e^{-x_a}}2=\sinh x_a.
 \label{eq:planck-ct108-lectores-par-impar-rev3}
\end{equation}
Le premier conserve la distance au point autodual et le second conserve son
orientation.

Les lecteurs physiques réciproques associés à la même coordonnée sont
\[
 r_{{\rm g},a}(m)=\frac{G_{108,a}^{\circ}m}{c_{\rm int}^{2}},
 \qquad
 \bar\lambda_{{\rm C},a}(m)=\frac{\hbar_a}{mc_{\rm int}},
\]
et satisfont
\begin{equation}
 \frac{r_{{\rm g},a}(m)}{\bar\lambda_{{\rm C},a}(m)}
 =e^{2x_a(m)},
 \qquad
 \frac{\bar\lambda_{{\rm C},a}(m)}{r_{{\rm g},a}(m)}
 =e^{-2x_a(m)}.
 \label{eq:planck-ct108-lectores-exponenciales-rev3}
\end{equation}
La conjugaison échange les deux lecteurs ; l’identité de Planck est leur
équilibre unitaire.

\begin{exactbox}[title={Identité de Planck et clôture électronique}]
L'identité \(m=m_{{\rm P},a}^{\circ}\) fixe l'origine autoduale de la droite
positive des masses. La section électronique fixe une position distincte sur
cette même droite au moyen d'une généalogie propre. La
\cref{sec:ley-general-masa-accion-autonoma} construit l'opérateur universel
qui transporte échelle absolue, caractère, tours et lecture bidirectionnelle ;
la \cref{sec:cierre-electron-two-way-rev11} réalise ensuite sa réduction
centrale de rang un. Ainsi, l'unité de Planck fournit l'identité de
la coordonnée et l'électron fournit la première transition matérielle
complète sur celle-ci.
\end{exactbox}

\begin{proposition}[Réalisation généalogique du centre électronique]
\label{prop:planck-ct108-realizacion-centro-electron-rev3}
Soit
\[
 F_e^+=\mathbb R_{>0}(P_++P_-)
\]
la demi-droite positive de la fibre centrale de rang un, soit
\(\mathcal G_e\) l’espace des généalogies
électroniques enrichies et soit \(\mathcal L_M^+\) la droite positive des
masses. Un morphisme de réalisation
\[
 \mathfrak R_e:
 F_e^+\times\mathcal G_e\longrightarrow\mathcal L_M^+
\]
peut envoyer l’identité \(I=P_++P_-\) sur une masse différente de
\(m_{{\rm P},a}^{\circ}\) sans modifier le
\cref{thm:planck-ct108-recta-partida-masa-rev3} : l’identité de Planck fixe
l’origine autoduale de la carte, tandis que la généalogie électronique fixe un
déplacement sur cette carte.
\end{proposition}

\begin{proof}
Si \(\chi_e:\mathcal G_e\to\mathbb R_{>0}\) est le caractère positif construit
par le registre électronique, la règle
\[
 \mathfrak R_e(\lambda I,g)
 =\lambda\,m_{{\rm P},a}^{\circ}\chi_e(g)
\]
est positive et multiplicative dans la coordonnée généalogique. Pour la généalogie
neutre, \(\chi_e=1\), on retrouve l'origine de Planck ; pour une généalogie avec
\(\chi_e(g)\ne1\), la même identité interne occupe une position distincte sur
la droite. La clôture électronique concrète conserve sa source dans la
\cref{sec:cierre-electron-two-way-rev11}.
\end{proof}

\paragraph{Domaine de l'identité autoduale de Planck : périodicité temporelle d'ordre \(108\), réalisation circulaire déclarée et séparation de la clôture électronique}
La droite d'action pré-retour/post-retour et la périodicité temporelle
d'ordre (108) précèdent la réalisation circulaire, qui est prise comme primitive
géométrique déclarée. L'identité entre la périodicité
circulaire et les lectures de Compton, le
sélecteur \(\theta_{108}^{\rm circ}\), l'unité interne de Planck et
l'involution de la droite positive sont des résultats exacts une fois cette
réalisation fixée. La sélection gravitationnelle générale et la clôture électronique
sont construites séparément dans les sections indiquées.
